% 2011 年安徽卷理科第 6 题：几何体的三视图（正视图 4×4 正方形、侧视图梯形、俯视图 4×4 分 1–2–1）。
% 来源 PDF：普通高考/2011/2011安徽理.pdf，SHA-256 83dabf28711e3020c00f9fc73f3bf7be70f65fed94ab6ad5b6a39ac840820c0f。
% 原页：第 1 页（A3 三栏）第 1 栏中下部，书页码 643。
% 重绘依据：照原页布局——正（主）视图在左上，侧（左）视图在右上，俯视图在左下；
%   正视图为边长 4 的正方形，左侧与上方各标 4；侧视图为梯形，下底与原页同宽、上底约为其一半、高 4；
%   俯视图为 4×4 矩形，自上而下由两条横线分成 1、2、1 三段，左侧标 1、2、1。
%   原页未在侧视图上标上底长度，此处不补注数字；按原页拓扑与相对比例绘，不按原页绝对比例。
\begin{tikzpicture}[scale=0.95, >=stealth, semithick, font=\small]
  % 正视图（4×4）
  \draw (0,0) rectangle (4,4);
  % 侧视图（梯形：下底 4，上底 2，高 4）
  \draw (6,0) -- (10,0) -- (9,4) -- (7,4) -- cycle;
  % 俯视图（4×4，自上而下 1、2、1）
  \draw (0,-1.5) rectangle (4,-5.5);
  \draw (0,-2.5) -- (4,-2.5);
  \draw (0,-4.5) -- (4,-4.5);
  % 正视图尺寸标注
  \draw[<->] (0,4.42) -- (4,4.42) node[midway,above] {$4$};
  \draw[<->] (-0.52,0) -- (-0.52,4) node[midway,left] {$4$};
  % 俯视图尺寸标注
  \draw[<->] (-0.52,-1.5) -- (-0.52,-2.5) node[midway,left] {$1$};
  \draw[<->] (-0.52,-2.5) -- (-0.52,-4.5) node[midway,left] {$2$};
  \draw[<->] (-0.52,-4.5) -- (-0.52,-5.5) node[midway,left] {$1$};
  % 视图名称
  \node at (2,-0.65) {正（主）视图};
  \node at (8,-0.65) {侧（左）视图};
  \node at (2,-6.1) {俯视图};
\end{tikzpicture}
